% =============================================================================
%  QuickBites — Assignment 4, SubTask 1
%  Shard Key Selection & Justification
% =============================================================================
\documentclass[12pt, a4paper]{article}

% ---- Packages ---------------------------------------------------------------
\usepackage[margin=2.5cm]{geometry}
\usepackage{graphicx}
\usepackage{booktabs}
\usepackage{longtable}
\usepackage{tabularx}
\usepackage{array}
\usepackage{hyperref}
\usepackage{listings}
\usepackage{xcolor}
\usepackage{fancyhdr}
\usepackage{caption}
\usepackage{amsmath}
\usepackage{enumitem}
\usepackage{parskip}
\usepackage{multirow}
\usepackage{amssymb}
\usepackage{colortbl}
\captionsetup[longtable]{justification=centering}

% ---- Table layout helpers ---------------------------------------------------
\setlength{\LTleft}{0pt}
\setlength{\LTright}{0pt}
\renewcommand{\arraystretch}{1.15}
\newcolumntype{L}[1]{>{\raggedright\arraybackslash}p{#1}}
\newcolumntype{C}[1]{>{\centering\arraybackslash}p{#1}}
\newcolumntype{R}[1]{>{\raggedleft\arraybackslash}p{#1}}
\newcolumntype{Y}{>{\raggedright\arraybackslash}X}
\let\plainunderscore\_
\renewcommand{\_}{\plainunderscore\allowbreak}

% ---- Colours ----------------------------------------------------------------
\definecolor{sqlkw}{RGB}{0,0,180}
\definecolor{sqlstr}{RGB}{163,21,21}
\definecolor{sqlcmt}{RGB}{0,128,0}
\definecolor{codebg}{RGB}{248,248,248}
\definecolor{HeaderGray}{HTML}{E6E6E6}
\definecolor{shard0}{HTML}{D6EAF8}
\definecolor{shard1}{HTML}{D5F5E3}
\definecolor{shard2}{HTML}{FDEBD0}

% ---- Code listing style -----------------------------------------------------
\lstdefinelanguage{SQL}{
  keywords={SELECT,FROM,WHERE,GROUP,BY,ORDER,UNION,ALL,AS,MOD,COUNT,MAX,MIN,AVG,
            SUM,ON,LEFT,JOIN,DISTINCT,USE,AND,OR,NOT,NULL,INNER},
  sensitive=false,
  morecomment=[l]{--},
  morestring=[b]',
}
\lstset{
    basicstyle=\ttfamily\small,
    keywordstyle=\color{sqlkw}\bfseries,
    stringstyle=\color{sqlstr},
    commentstyle=\color{sqlcmt}\itshape,
    backgroundcolor=\color{codebg},
    frame=single,
    rulecolor=\color{black!20},
    breaklines=true,
    showstringspaces=false,
    tabsize=4,
    captionpos=b,
    numbers=none,
    xleftmargin=4pt,
    xrightmargin=4pt,
    aboveskip=8pt,
    belowskip=8pt,
}
\lstdefinestyle{python}{language=Python,
    keywordstyle=\color{sqlkw}\bfseries,
    stringstyle=\color{sqlstr},
    commentstyle=\color{sqlcmt}\itshape,
}

% ---- Headers & footers ------------------------------------------------------
\pagestyle{fancy}
\fancyhf{}
\fancyhead[L]{\small QuickBites --- Assignment 4: Sharding}
\fancyhead[R]{\small SubTask 1: Shard Key Selection}
\fancyfoot[C]{\thepage}
\renewcommand{\headrulewidth}{0.4pt}

% ---- Hyperref setup ---------------------------------------------------------
\hypersetup{
    colorlinks=true,
    linkcolor=blue!60!black,
    citecolor=blue!60!black,
    urlcolor=blue!60!black,
}

% ---- Cover-page metadata ----------------------------------------------------
\newcommand{\projectname}{QuickBites}
\newcommand{\projectsubtitle}{A Food Delivery Platform}
\newcommand{\reportfocus}{Assignment~4 --- Sharding of the Developed Application}
\newcommand{\subtaskfocus}{SubTask~1: Shard Key Selection \& Justification}
\newcommand{\coursecode}{CS 432: Databases}
\newcommand{\assignmentinfo}{Assignment 4}
\newcommand{\semesterinfo}{Semester II, 2025--2026}
\newcommand{\repositoryurl}{https://github.com/mohit-codes-21/QuickBites}
\newcommand{\institutionname}{Indian Institute of Technology, Gandhinagar}

% =============================================================================
\begin{document}
\begin{titlepage}
\thispagestyle{empty}
\begin{center}
\vspace*{1.2cm}

{\Huge \textbf{\projectname}}\\[8pt]
{\LARGE \projectsubtitle}\\[12pt]
{\Large \reportfocus}\\[6pt]
{\large \subtaskfocus}\\[42pt]

{\Large \textbf{\coursecode}}\\[4pt]
{\Large \assignmentinfo}\\[28pt]

{\Large \semesterinfo}

\vfill

\begin{minipage}{0.84\textwidth}
\centering
\begin{tabular}{@{} L{5.0cm} L{7.0cm} @{}}
\toprule
\textbf{Team Members} & \textbf{Role} \\
\midrule
Bhavya Parmar        & Backend \& Database Design \\
Surriya Gokul        & SQL Indexing \& Benchmarking \\
Sonagiri Thoshith Kautilya & Frontend \& UI/UX \\
\bottomrule
\end{tabular}
\end{minipage}

\vspace{1.2cm}
{\large \textbf{GitHub:}
  \href{\repositoryurl}{\texttt{github.com/mohit-codes-21/QuickBites}}}

\vfill

{\Large \institutionname}
\end{center}
\end{titlepage}

\clearpage
\setcounter{page}{1}
\tableofcontents
\newpage

% =============================================================================
\section{Introduction}
% =============================================================================

This report documents \textbf{SubTask~1} of Assignment~4 for the QuickBites food
delivery platform. The objective of SubTask~1 is to:

\begin{enumerate}[leftmargin=*, label=\arabic*.]
  \item Select a \emph{Shard Key} that satisfies the three assignment criteria
        (high cardinality, query-aligned, and stable).
  \item Identify and justify the \emph{partitioning strategy} to be used.
  \item Estimate the \emph{expected data distribution} across shards and
        identify any risk of skew.
\end{enumerate}

SubTask~1 lays the analytical and design foundation on which SubTasks~2
(data partitioning), 3 (query routing), and 4 (scalability analysis) will
build. No shard tables are created yet; this document is the design artefact.

% =============================================================================
\section{QuickBites Schema Overview}
% =============================================================================

The QuickBites database (established in Assignment~1) consists of the thirteen
relations shown in Table~\ref{tab:schema}. The schema models a food delivery
marketplace with customers, restaurants, menus, orders, payments, and
delivery partners.

\begin{table}[h!]
\centering
\caption{QuickBites relational schema summary}
\label{tab:schema}
\begin{tabular}{@{} L{3.6cm} L{6.0cm} R{2.0cm} @{}}
\toprule
\textbf{Table} & \textbf{Key Columns} & \textbf{Rows (seed)} \\
\midrule
Member              & memberID (PK), name, email           & 30 \\
Customer            & customerID (PK, FK$\to$Member)       & 10 \\
DeliveryPartner     & partnerID (PK, FK$\to$Member)        & 10 \\
Restaurant          & restaurantID (PK), city, zipCode     & 10 \\
MenuItem            & (restaurantID, itemID) PK            & 20 \\
Address             & (customerID, addressID) PK, city     & 20 \\
CartItem            & (customerID, restaurantID, itemID)   & 12 \\
Payment             & paymentID (PK), customerID           & 12 \\
Orders              & orderID (PK), customerID, restaurantID & 12 \\
OrderItem           & (orderID, restaurantID, itemID)      & 22 \\
Delivery\_Assignments & AssignmentID (PK), OrderID, PartnerID & 12 \\
OrderRating         & orderID (PK)                         & 12 \\
MenuItemRating      & (restaurantID, itemID, orderID)      & 22 \\
\bottomrule
\end{tabular}
\end{table}

% =============================================================================
\section{Shard Key Candidates}
% =============================================================================

To choose the best shard key we first surveyed all columns that
(a)~appear frequently in \texttt{WHERE} clauses of the Assignment~2 APIs, and
(b)~have enough distinct values to spread rows evenly.
Table~\ref{tab:candidates} lists the shortlisted candidates.

\begin{table}[h!]
\centering
\caption{Shard key candidate evaluation}
\label{tab:candidates}
\begin{tabular}{@{} L{3.2cm} C{2.2cm} C{2.4cm} C{2.0cm} C{1.8cm} @{}}
\toprule
\textbf{Candidate} & \textbf{High Cardinality} & \textbf{Query-Aligned} & \textbf{Stable} & \textbf{Score} \\
\midrule
\textbf{customerID} & \checkmark & \checkmark\checkmark & \checkmark & \textbf{5/5} \\
restaurantID        & \checkmark & \checkmark           & \checkmark & 4/5 \\
city (Restaurant)   & \texttimes\ (low) & \checkmark   & \checkmark & 3/5 \\
zipCode             & \checkmark & \texttimes           & \checkmark & 3/5 \\
orderID (hash)      & \checkmark & \checkmark           & \checkmark & 4/5 \\
\bottomrule
\end{tabular}
\end{table}

\noindent\textbf{Selected shard key: \texttt{customerID}}

% =============================================================================
\section{Shard Key Justification}
% =============================================================================

\subsection{Criterion 1 --- High Cardinality}

\texttt{customerID} is a unique integer auto-assigned at registration.
With $N$ registered customers there are $N$ distinct values, so the hash
function distributes rows over $k$ shards with an expected imbalance of
$O(1/\sqrt{N})$ --- negligible for production scale.

In the seed dataset, the cardinality survey (Listing~\ref{lst:cardinality})
confirms:

\begin{center}
\begin{tabular}{lrr}
\toprule
Candidate key & Distinct values & Total rows \\
\midrule
customerID          & 10 & 10 \\
restaurantID        & 10 & 10 \\
city (Restaurant)   &  4 & 10 \\
zipCode             &  7 & 10 \\
\bottomrule
\end{tabular}
\end{center}

\noindent \texttt{city} has only 4 distinct values, making it unsuitable; the
remaining candidates are equivalent on cardinality. The decision falls on
the next two criteria.

\subsection{Criterion 2 --- Query-Aligned}

A key is \emph{query-aligned} if it appears in the \texttt{WHERE} clause of
the most frequent API calls. Analysing the Assignment~2 Flask endpoints:

\begin{itemize}[leftmargin=*]
  \item \textbf{Place order / view cart} --- keyed by \texttt{customerID}
  \item \textbf{Payment history} --- filtered by \texttt{customerID}
  \item \textbf{Address book (CRUD)} --- filtered by \texttt{customerID}
  \item \textbf{Order status / tracking} --- \texttt{orderID} which resolves
        to a \texttt{customerID} via the \texttt{Orders} table
  \item \textbf{Order rating submission} --- \texttt{orderID}
        $\to$ \texttt{customerID}
\end{itemize}

Routing on \texttt{customerID} means that the vast majority of these
operations touch \emph{exactly one shard}, eliminating cross-shard joins
for the hot path.

By contrast, routing on \texttt{restaurantID} would scatter a single
customer's order history across all shards, forcing scatter--gather for
every customer-facing read.

\subsection{Criterion 3 --- Stability}

\texttt{customerID} is assigned once at account creation and is never
updated. This satisfies the stability criterion: a customer's data is
written to shard $s = \texttt{customerID} \bmod k$ once and never needs
to be migrated to a different shard due to a key-value change.

In contrast, a mutable attribute such as \texttt{city} (delivery address)
would require data migration whenever a customer moves.

% =============================================================================
\section{Partitioning Strategy}
% =============================================================================

\subsection{Strategy Chosen: Hash-Based Modulo}

\[
  \text{shard\_id} \;=\; \texttt{customerID} \bmod \texttt{NUM\_SHARDS}
\]

with $\texttt{NUM\_SHARDS} = 3$.

\subsection{Alternatives Considered}

\begin{description}[leftmargin=2.2cm, style=nextline]
  \item[Range-based]
    Divide customerIDs into contiguous ranges
    (e.g.\ 1--1000 $\to$ Shard~0, 1001--2000 $\to$ Shard~1).
    \emph{Rejected}: new customers arrive with monotonically increasing IDs,
    so all new writes would land on the last shard, creating a write hotspot.

  \item[Hash-based (chosen)]
    Apply a deterministic hash to the shard key.
    \emph{Accepted}: uniformly spreads writes and reads; no hotspots under
    sequential ID assignment; O(1) routing with no external lookup.

  \item[Directory-based]
    Maintain a lookup table mapping key ranges to shard IDs.
    \emph{Rejected}: adds a round-trip to the directory on every request,
    introducing latency and a single point of failure. Suitable when shards
    must be remapped non-uniformly, which is not a current requirement.
\end{description}

\subsection{Why Hash-Based Suits QuickBites}

\begin{enumerate}[leftmargin=*, label=\arabic*.]
  \item \textbf{Even write distribution.} Customer registrations are
        sequential, but modulo hashing cycles over all shards, so no shard
        receives all new writes.
  \item \textbf{Deterministic and stateless routing.} The routing function
        is a single arithmetic operation; the application layer requires no
        extra database call to determine which shard to contact.
  \item \textbf{Consistent read locality.} All rows for a given
        \texttt{customerID} (orders, payments, addresses, cart items) hash
        to the same shard, so customer-scoped queries remain single-shard.
  \item \textbf{Simplicity.} The formula is trivially expressible in both
        Python and SQL.
\end{enumerate}

% =============================================================================
\section{Expected Data Distribution}
% =============================================================================

\subsection{Routing Formula Applied to Seed Data}

Table~\ref{tab:assignment} shows the shard assignment for each of the ten
seed customers.

\begin{table}[h!]
\centering
\caption{CustomerID $\to$ Shard assignment (\texttt{customerID MOD 3})}
\label{tab:assignment}
\begin{tabular}{@{} C{2.5cm} C{2.5cm} C{4.0cm} @{}}
\toprule
\textbf{customerID} & \textbf{Shard ID} & \textbf{CustomerIDs in shard} \\
\midrule
\rowcolor{shard0} 3, 6, 9     & 0 & \{3, 6, 9\} \\
\rowcolor{shard1} 1, 4, 7, 10 & 1 & \{1, 4, 7, 10\} \\
\rowcolor{shard2} 2, 5, 8     & 2 & \{2, 5, 8\} \\
\bottomrule
\end{tabular}
\end{table}

\subsection{Row-Count Distribution Across Tables}

Table~\ref{tab:distribution} shows the expected number of rows per shard
for every customer-linked table in the seed dataset. Tables that are
\emph{not} linked to \texttt{customerID} (Restaurant, MenuItem,
DeliveryPartner, Member, Delivery\_Assignments, OrderRating,
MenuItemRating) are classified as \emph{global/reference} tables and will
be replicated across all shards rather than partitioned.

\begin{table}[h!]
\centering
\caption{Expected row distribution across three shards (seed dataset)}
\label{tab:distribution}
\begin{tabular}{@{} L{3.0cm} C{2.0cm} C{2.0cm} C{2.0cm} C{2.0cm} C{1.8cm} @{}}
\toprule
\textbf{Table} & \textbf{Shard~0} & \textbf{Shard~1} & \textbf{Shard~2} & \textbf{Total} & \textbf{Max/Avg} \\
\midrule
Customer    &  3 &  4 &  3 & 10 & 1.20 \\
Orders      &  4 &  5 &  3 & 12 & 1.25 \\
Payment     &  4 &  5 &  3 & 12 & 1.25 \\
Address     &  6 &  8 &  6 & 20 & 1.20 \\
CartItem    &  4 &  5 &  3 & 12 & 1.25 \\
\midrule
\textbf{Total} & \textbf{21} & \textbf{27} & \textbf{18} & \textbf{66} & \textbf{1.23} \\
\bottomrule
\end{tabular}
\end{table}

\noindent The \emph{skew coefficient} (maximum shard rows / average shard rows)
for each table is at most $1.25$, well within the acceptable range ($< 1.5$).

\subsection{Verification: SQL Query for Distribution}

Listing~\ref{lst:distribution} shows the SQL query (from
\texttt{subtask1\_shard\_analysis.sql}) that replicates the distribution
table above directly against the live \texttt{QB} database.

\begin{lstlisting}[language=SQL, caption={Distribution verification query},
                   label={lst:distribution}]
SELECT
    shard_id,
    SUM(customer_rows) AS customers,
    SUM(order_rows)    AS orders,
    SUM(payment_rows)  AS payments,
    SUM(address_rows)  AS addresses,
    SUM(cart_item_rows) AS cart_items
FROM (
    SELECT customerID MOD 3 AS shard_id,
           1 AS customer_rows, 0, 0, 0, 0
    FROM Customer
    UNION ALL
    SELECT customerID MOD 3, 0, 1, 0, 0, 0 FROM Orders
    UNION ALL
    SELECT customerID MOD 3, 0, 0, 1, 0, 0 FROM Payment
    UNION ALL
    SELECT customerID MOD 3, 0, 0, 0, 1, 0 FROM Address
    UNION ALL
    SELECT customerID MOD 3, 0, 0, 0, 0, 1 FROM CartItem
) sub
GROUP BY shard_id
ORDER BY shard_id;
\end{lstlisting}

\subsection{Shard Routing Function}

Listing~\ref{lst:router} shows the shard routing helper implemented in
\texttt{A4/shard\_key\_analysis.py}. This function is the single
authoritative routing rule for the application.

\begin{lstlisting}[style=python, caption={Hash-based shard router},
                   label={lst:router}]
NUM_SHARDS = 3

def get_shard(customer_id: int, num_shards: int = NUM_SHARDS) -> int:
    """Return the shard index for a given customerID.

    Uses modulo hashing so that every customerID maps deterministically
    to exactly one shard without any external lookup table.
    """
    return customer_id % num_shards
\end{lstlisting}

\subsection{Cardinality Query}

Listing~\ref{lst:cardinality} shows the SQL query used to compare
cardinalities across all shard key candidates.

\begin{lstlisting}[language=SQL, caption={Candidate key cardinality survey},
                   label={lst:cardinality}]
SELECT 'customerID' AS candidate_key,
       COUNT(DISTINCT customerID) AS distinct_values
FROM Customer
UNION ALL
SELECT 'restaurantID', COUNT(DISTINCT restaurantID) FROM Restaurant
UNION ALL
SELECT 'city (Restaurant)', COUNT(DISTINCT city)    FROM Restaurant
UNION ALL
SELECT 'zipCode', COUNT(DISTINCT zipCode)           FROM Restaurant;
\end{lstlisting}

% =============================================================================
\section{Skew Risk Analysis}
% =============================================================================

\subsection{Current Skew (Seed Dataset)}

With 10 customers and 3 shards the modulo remainder produces a $4$--$3$--$3$
split. The skew coefficient is:

\[
  \text{skew} = \frac{\max(\{4,3,3\})}{\bar{x}} = \frac{4}{10/3} \approx 1.20
\]

A coefficient below $1.25$ is classified as \textbf{low skew} and is
acceptable for production deployment.

\subsection{Scalability of the Hash Distribution}

At larger scales the imbalance shrinks. For $N$ customers uniformly
distributed modulo $k$ shards, the expected maximum count is:

\[
  \mathbb{E}[\max\text{-shard}] = \frac{N}{k} + O\!\left(\sqrt{\frac{N\ln k}{k}}\right)
\]

The skew coefficient therefore converges to $1$ as $N \to \infty$,
confirming that hash-based modulo is appropriate for production scale.

\subsection{Identified Skew Risks}

\begin{enumerate}[leftmargin=*, label=\arabic*.]
  \item \textbf{Small-dataset remainder imbalance.}
        With only 10 customers and 3 shards, Shard~1 holds one extra customer
        (40\% vs.\ the ideal 33.3\%). This is an inherent consequence of
        $10 \not\equiv 0 \pmod{3}$ and disappears at scale.

  \item \textbf{Batch imports with correlated IDs.}
        If a bulk import assigns all even IDs to a single batch, those IDs
        will skew toward shards $0$ and $2$ only
        ($\text{even} \bmod 3 \in \{0, 2\}$). Mitigation: use a
        cryptographic hash (e.g.\ MurmurHash3) instead of plain modulo
        for large-scale deployment.

  \item \textbf{High-activity customers.}
        A single customer placing many orders does \emph{not} create hot
        shards in the sense of uneven row distribution, because all of their
        orders still go to the same shard. However, it may create a
        \emph{read/write hotspot} on that shard if one customer dominates
        query traffic. Mitigation: monitor per-shard QPS and rebalance by
        increasing shard count if needed.

  \item \textbf{Shard-local cross-table joins.}
        OrderItem, Delivery\_Assignments, OrderRating, and MenuItemRating
        are linked to Orders rather than directly to customerID. These will
        be co-located on the same shard as their parent Order (since
        Orders are routed by customerID), so most joins remain shard-local.
        Full cross-shard aggregation (e.g.\ global revenue reports) will
        require scatter--gather queries across all shards.
\end{enumerate}

% =============================================================================
\section{Global vs.\ Sharded Table Classification}
% =============================================================================

Not all tables benefit from sharding. Table~\ref{tab:classification}
classifies each QuickBites relation.

\begin{table}[h!]
\centering
\caption{Table sharding classification}
\label{tab:classification}
\begin{tabular}{@{} L{4.0cm} C{2.6cm} L{5.8cm} @{}}
\toprule
\textbf{Table} & \textbf{Classification} & \textbf{Rationale} \\
\midrule
Customer         & Sharded    & Primary shard-key owner \\
Orders           & Sharded    & FK on customerID \\
Payment          & Sharded    & FK on customerID \\
Address          & Sharded    & FK on customerID \\
CartItem         & Sharded    & FK on customerID \\
OrderItem        & Sharded    & FK on Orders (co-located) \\
\midrule
Restaurant       & Global/replicated & No customerID FK; needed by all shards \\
MenuItem         & Global/replicated & FK on Restaurant only \\
Member           & Global/replicated & Parent of Customer \& DeliveryPartner \\
DeliveryPartner  & Global/replicated & Assigned per order, not per customer \\
Delivery\_Assignments & Sharded (via Orders) & FK on OrderID (Order is sharded) \\
OrderRating      & Sharded (via Orders) & FK on orderID \\
MenuItemRating   & Sharded (via Orders) & FK on orderID \\
\bottomrule
\end{tabular}
\end{table}

% =============================================================================
\section{Summary}
% =============================================================================

\begin{itemize}[leftmargin=*]
  \item \textbf{Shard key selected:} \texttt{customerID}
  \item \textbf{Criteria satisfied:} High cardinality ($N$ distinct values
        for $N$ customers), query-aligned (present in $>$90\% of API
        WHERE clauses), and stable (never updated after insertion).
  \item \textbf{Partitioning strategy:} Hash-based modulo
        ($\texttt{customerID} \bmod 3$), chosen for even distribution,
        deterministic O(1) routing, and no external lookup overhead.
  \item \textbf{Expected distribution:} Shard~0 $\approx 30\%$,
        Shard~1 $\approx 40\%$, Shard~2 $\approx 30\%$ (seed dataset).
        Skew coefficient $\approx 1.20$ (low skew).
  \item \textbf{Skew risks identified:} Small-dataset remainder imbalance,
        batch imports with correlated IDs, and per-customer activity
        hotspots --- all of which are manageable and noted for SubTask~4.
\end{itemize}

The design is implemented in \texttt{A4/shard\_key\_analysis.py}
(Python routing function and distribution analysis) and
\texttt{A4/subtask1\_shard\_analysis.sql} (SQL distribution
verification queries). These artefacts provide the foundation for
SubTask~2 (shard table creation and data migration) and SubTask~3
(query routing implementation).

% =============================================================================
\appendix
\section{Repository Structure for Assignment 4}
% =============================================================================

\begin{verbatim}
A4/
+-- Track1_assignment4.pdf          # Assignment spec
+-- A4_Subtask1_Report.tex          # This report (SubTask 1)
+-- shard_key_analysis.py           # Routing function + distribution analysis
+-- subtask1_shard_analysis.sql     # SQL cardinality + distribution queries
\end{verbatim}

\end{document}
